\documentclass[manuscript,screen,review]{acmart}

% --- CONFIGURAÇÕES DE LOGO DO BIBTEX ---
\AtBeginDocument{%
	\providecommand\BibTeX{{%
			Bib\TeX}}}

% --- INFORMAÇÕES DE DIREITOS AUTORAIS ---
% IMPORTANTE: Deixado como 'none' para a fase de escrita.
% A ACM te enviará os códigos corretos após a aceitação do artigo.
\setcopyright{none}
\copyrightyear{2026}
\acmYear{2026}
\acmDOI{}
\acmConference[Conferência '26]{Nome da Conferência}{Junho 2026}{Local}
\acmISBN{}

% --- METADADOS DO DOCUMENTO ---
\title[Análise de Desempenho de Modelos de Linguagem]{Análise de Desempenho da Inferência de um Modelo de Linguagem Particionado em Múltiplas GPUs}

% Autor 1
\author{Lucas Fraga Balbinot}
\email{fraga.balbinot@ufrgs.br}
\affiliation{%
	\institution{Universidade Federal do Rio Grande do Sul (UFRGS)}
	\city{Porto Alegre}
	\state{Rio Grande do Sul}
	\country{Brasil}
}

% Autor 2 (Se houver mais autores, basta copiar e colar este bloco abaixo)
\author{Matheus Augusto Tregnago}
\email{autor2@instituicao.edu.br}
\affiliation{%
	\institution{Universidade Federal do Rio Grande do Sul (UFRGS)}
	\city{Porto Alegre}
	\state{Rio Grande do Sul}
	\country{Brasil}
}

\author{Rafael Silva de Souza}
\email{autor2@instituicao.edu.br}
\affiliation{%
	\institution{Universidade Federal do Rio Grande do Sul (UFRGS)}
	\city{Porto Alegre}
	\state{Rio Grande do Sul}
	\country{Brasil}
}

% Define como os autores aparecem no cabeçalho das páginas ímpares
\renewcommand{\shortauthors}{Balbinot, Tregnago, Souza}

% --- CORPO DO DOCUMENTO ---
\begin{document}

	% O Resumo (Abstract) DEVE vir antes do \maketitle
	\begin{abstract}
		Estudo do comportamento de Modelos de Linguagem de Grande Porte (\textit{Large Language Models}, ou LLM) em situação de particionamento de dados, utilizando as GPUs nativas aos sistemas da Universidade Federal do Rio Grande do Sul, Campus do Vale.
	\end{abstract}

	% --- PALAVRAS-CHAVE ---
	\keywords{Large Language Model, Particionamento de Dados, GPU, Análise de Desempenho}

	\maketitle

	% --- SEÇÕES DO ARTIGO ---
	\section{Introdução}
	Na era moderna, a prevalência de Inteligências Artificiais (\textit{Artificial Intelligences} em inglês, abreviação AI) no âmbito diário é cada vez mais notável, seja na forma de assistentes virtuais, de sistemas de tradução ou mesmo na geração de códigos para programas. Um elemento essencial para sua aplicabilidade cada vez mais abrangente são os Modelos de Linguagem de Grande Porte (\textit{Large Language Models} ou LLM), vastas redes neurais capazes de processar bilhões de parâmetros dados por vários usuários eficientemente e tendo tremendo poder computacional para tal tarefa. Contudo, tais atributos também fazem com que LLMs ocupem muito espaço em máquinas, por vezes necessitando que seja particionada em segmentos menores para poder operar.

	O propósito deste trabalho é o estudo dos efeitos de particionamento de LLMs por Unidades de Processamento Gráfico (\textit{Graphic Process Unit} ou GPU) e o estudo de como afeta o processo de inferência do modelo, utilizando as GPUs localizadas no Parque Computacional de Alto Desempenho (PCAD) da Universidade Federal do Rio Grande do Sul (UFRGS), e mantida pelo Laboratório de Processamento Paralelo e Distribuído (LPPD), como ambiente de teste.

	\section{Contexto e Trabalhos Relacionados}
	Um problema fundamental no desdobramento de LLMs reside na alta demanda de memória necessária para sua operação. Grande parte desses modelos fundamenta-se na arquitetura Transformer \cite{VASWANI2017}, a qual utiliza mecanismos de atenção para modelar relações de dependência em dados sequenciais com alta precisão. Contudo, mesmo as arquiteturas mais modernas enfrentam severas limitações decorrentes do custo computacional tanto no treinamento quanto na inferência \cite{BOMMASANI2021}, o que gera uma dependência estrita de clusters de Unidades de Processamento Gráfico (GPUs) para viabilizar sua execução.

	Embora o paralelismo de modelos seja amplamente investigado na literatura recente, a análise do impacto do particionamento de inferência em ambientes acadêmicos compartilhados — como o PCAD da UFRGS — ainda apresenta lacunas que este trabalho busca explorar.

	Um ambiente capaz de computação de alto desempenho é vital para este estudo. Nesse sentido, o PCAD prova-se ideal, visto que possui uma grande variedade de modelos que levam a exploração de diferentes reações de particionamento com requisitos heterogêneos.

	\section{Metodologia Proposta}
	Descrição do método, algoritmos ou arquitetura proposta.

	O modelo utilizado no estudo foi o Qwen2.5-7B-Instruct, um modelo usado para pré e pós-treinamento de IA, possuindo 7.61 bilhões de parâmetros no total e 28 camadas\cite{qwen2.5_7b_instruct}. Tais qualidades foram consideradas satisfatórias para o propósito do estudo, dando

	Para poder estudar o comportamento das GPUs, foram definidas três configurações de particionamento diferentes:

	\begin{enumerate}
	\item 1 GPU;
	\item 2 GPUs;
	\item 4 GPUs.
	\end{enumerate}

	Para as configurações 2 e 3, há uma complexidade adicional em que elas funcionam segundo dois comportamentos de paralelismo, com o comportamento de GPU único funcionando como o \textit{baseline} dos sistemas. Os tipos de paralelismo são:

	\begin{enumerate}
		\item \textit{Tensor Parallelism}: Paralelismo baseado na fragmentação das camadas de um modelo em diferentes máquinas;
		\item \textit{Pipeline Parallelism}: Paralelismo em que a divisão das camadas é mais processado como um "estado", passando cada instrução sequencialmente entre cada máquina.
	\end{enumerate}

	Isso foi feito para poder fazer uma análise objetiva do efeito que o \textit{overhead} causado por paralelismo causa na execução de um programa, e como a capacidade de comunicação entre máquinas afeta seu desempenho como um todo.

	Dos atributos analisados, aqueles que foram dados maior atenção em particular são os que relatam aos tokens: ou seja, o tempo total até o primeiro token (referido como TTFT); o tempo médio entre tokens consecutivos (ITL); a latência de pedido; e o throughput de tokens. Ademais, foi analisado também o comportamento da GPU nessas circunstâncias, focando na banda de memória usada (\textit{memory bandwidth}), o quanto da memória é usada (\textit{GPU memory usage}) e a ocupação dos multiprocessadores de streaming (\textit{SM Occupancy}). Para o \textit{benchmarking} destas métricas,

	\section{Resultados e Avaliação}
	Resultados experimentais obtidos e discussões.

	\section{Conclusão}
	Considerações finais e propostas de trabalhos futuros.

	% --- AGRADECIMENTOS ---
	\begin{acks}
		Agradecemos ao LPPD da UFRGS por poder providenciar o espaço necessário para realizar estes experimentos, além da variedade oferecida.
		Agradecemos ao Professor Lucas Mello Schnorr pela orientação oferecida no desenvolvimento deste trabalho e seus conselhos.
\end{acks}

	% --- BIBLIOGRAFIA ---
	\bibliographystyle{ACM-Reference-Format}
	\bibliography{acmart.bib} % Certifique-se de que o arquivo 'sample-base.bib' existe na pasta

\end{document}
\endinput
